\documentclass[11pt]{article}
\usepackage[T1]{fontenc}
\usepackage[utf8]{inputenc}
\usepackage{lmodern,microtype}
\usepackage[margin=0.88in]{geometry}
\usepackage{amsmath,amssymb,amsthm,mathtools}
\usepackage{booktabs,longtable,array}
\usepackage{xcolor,fancyhdr,enumitem,listings}
\usepackage[colorlinks=true,linkcolor=blue!40!black,citecolor=blue!40!black,urlcolor=blue!40!black]{hyperref}
\usepackage{url}
\hypersetup{pdftitle={Strict Horizons and a Free Quotient: Smaller Explicit Diophantine Equations for ZFC Consistency},pdfauthor={AI-assisted research prepared for Vladimir Reshetnikov and ProveIt},pdfsubject={Explicit Diophantine polynomials, arithmetic proofs, and reproducible finite certificates}}
\definecolor{ink}{RGB}{25,44,65}
\pagestyle{fancy}
\fancyhf{}
\fancyhead[L]{\small Strict horizons and a free quotient}
\fancyhead[R]{\small ProveIt research contribution}
\fancyfoot[C]{\thepage}
\setlength{\headheight}{14pt}
\setlength{\parindent}{0pt}
\setlength{\parskip}{0.35em}
\setlist{nosep,leftmargin=*}
\allowdisplaybreaks
\emergencystretch=2em
\newtheorem{theorem}{Theorem}[section]
\newtheorem{lemma}[theorem]{Lemma}
\newtheorem{proposition}[theorem]{Proposition}
\newtheorem{corollary}[theorem]{Corollary}
\theoremstyle{definition}
\newtheorem{definition}[theorem]{Definition}
\newtheorem{remark}[theorem]{Remark}
\newcommand{\Z}{\mathbb Z}
\newcommand{\N}{\mathbb N}
\newcommand{\Con}{\operatorname{Con}}
\newcommand{\lcm}{\operatorname{lcm}}
\newcommand{\LL}{\mathcal L}
\newcommand{\calC}{C_{222}}
\lstset{basicstyle=\ttfamily\small,breaklines=true,columns=fullflexible,keepspaces=true,frame=single,rulecolor=\color{black!20}}
\title{\color{ink}\bfseries Strict Horizons and a Free Quotient\par
\large Smaller explicit Diophantine equations for ZFC consistency}
\author{AI-assisted mathematical research\smallskip\\
\normalsize Prepared for Vladimir Reshetnikov and the ProveIt project}
\date{25 October 2026}

\begin{document}
\maketitle
\thispagestyle{empty}
\begin{abstract}
We give a fully specified quartic polynomial in 236 unrestricted integer
variables, with 4,876 collected monomials and coefficient height 7,968,
whose absence of integer zeros is provably equivalent in ZFC to
$\Con(\mathrm{ZFC})$. It is the sum of the squares of 173 explicitly
printed quadratic polynomials, each of coefficient height at most 64.
A denser companion has 235 variables. The direct construction uses the
repository's numerical proof calculus and contains no large machine
coefficient or specialization of a universal Diophantine polynomial.

The principal new step chooses the proof horizon \emph{strictly} above
all local schema parameters and derived labels. Five range certificates
then share an existing binomial denominator, removing an entire exact
power request and its auxiliary data. A self-power majorant $F=t^t$,
a linear CRT radius $F$, and an affine signed radix $q=4z+1$ give further
compatible savings. Both directions are proved, including sign recovery
before binomial decoding and simultaneous cofinal witness choices.

A separate compact construction moves an arithmetic scale from a quotient
coordinate into its Pell base. It gives ten-positive-variable polynomials
with 83 binary operations and degree 312, or 89 operations and degree 278.
These retain the complete large fixed coefficient tuple and its inherited
machine calibration. Exact algebraic audits and reconstructors accompany
all four equations. The new unbounded arguments are ordinary proofs;
the inherited logical and Pell results are identified, and no new complete
proof-assistant formalization or global minimality claim is made.
\end{abstract}
\vfill
\begin{center}\small
Repository comparison snapshot:\\
\texttt{7a8da414376a80b90ecf799f37ceb26fb9ac674e}.\\
The equations are fixed; no consistency assumption is used to prove their calibration.
\end{center}
\newpage
\tableofcontents
\newpage

\section{The result and the choice of objective}
\label{sec:result}

The recommended small-coefficient equation is
\begin{equation}\label{eq:principal}
\boxed{\quad P_{236}(x_0,\ldots,x_{235})
   =\sum_{j=0}^{172}R_j(x_0,\ldots,x_{235})^2=0,
   \qquad \boldsymbol x\in\Z^{236}.\quad}
\end{equation}
Every $R_j$ is printed in Appendix~\ref{sec:residuals}. This is a literal
definition by small integer coefficients, not a prescription to compile
an unspecified machine. The archive contains the same residuals in
\nolinkurl{direct/principal236/equation.json}, their complete collected
quartic, readable exports, and exact transformation certificates.

\begin{theorem}[Fixed consistency equations]\label{thm:main}
For the explicitly supplied polynomials $P_{236}$ and $P_{235}$,
\[
\mathrm{ZFC}\vdash
\left(\neg\exists\boldsymbol x\in\Z^{236}\ P_{236}(\boldsymbol x)=0\right)
\longleftrightarrow\Con(\mathrm{ZFC}),
\]
and the corresponding statement holds with $235$ variables and $P_{235}$.
The two compact polynomials of Section~\ref{sec:free-quotient} satisfy
the analogous equivalence over ten strictly positive integers.
\end{theorem}

Here $\Con(\mathrm{ZFC})$ means the arithmetic statement that no finite
syntactic ZFC refutation exists, with the repository's specified proof
system. The proofs establish an equivalence of \emph{existence of
certificates}, without assuming consistency. The theorem does not assert
that ZFC proves its own consistency or the actual unsolvability of these
equations. It is an ordinary mathematical theorem obtained by composing
the new arithmetic proofs with the identified inherited calibration.
The verification programs check finite algebraic data, not this entire
unbounded theorem.

\subsection{Exact improvements over matched incoming equations}

\begin{center}\small
\begin{tabular}{lrrrrr}
\toprule
Unrestricted-integer quartic & Variables & Rows & Terms & Height & Row height\\
\midrule
Factorial-free parent, principal &242&178&5,018&7,968&64\\
\textbf{New principal} &\textbf{236}&\textbf{173}&\textbf{4,876}&\textbf{7,968}&\textbf{64}\\
Factorial-free parent, dense &241&177&12,774&7,968&64\\
\textbf{New dense companion} &\textbf{235}&\textbf{172}&\textbf{12,566}&\textbf{7,968}&\textbf{64}\\
\bottomrule
\end{tabular}
\end{center}

All four rows have exact total degree four. Thus each new equation saves
six variables and five residuals against its matched parent, while also
reducing support without increasing either coefficient height. The
principal has 142 fewer collected monomials; the dense companion has 208
fewer. The denser equation is a different tradeoff from the principal:
one coordinate is removed at a substantial increase in support.

\begin{center}\small
\begin{tabular}{lrrr}
\toprule
Ten-positive-variable circuit & Operations & Exact degree & Comparison\\
\midrule
Incoming plain point &84&328&retained predecessor\\
\textbf{New plain point} &\textbf{83}&\textbf{312}&both improve\\
Incoming corrected point &89&288&retained predecessor\\
Incoming corrected point &92&278&retained predecessor\\
\textbf{New corrected point} &\textbf{89}&\textbf{278}&dominates these two\\
\bottomrule
\end{tabular}
\end{center}

An operation is one binary addition, subtraction, or multiplication in
the live straight-line circuit. Reusing a computed value does not cost
another operation. The compact equations use a large fixed coefficient
tuple, whereas the direct quartics use only the displayed small
coefficients. These metrics must not be conflated. The repository already
contains equations with fewer variables and other favorable tradeoffs
\cite{repo}. We have neither a definition-independent smallest equation
nor a lower bound proving any of these endpoints optimal.

\subsection{What is new, and what is inherited}

The direct arithmetic contribution is the combination of a strict proof
horizon, a shared range denominator, self-power growth, a smaller CRT
radius, and a signed affine radix with a noncircular sign proof. The
compact contribution is a fresh placement of the quotient scale in the
Pell construction. These are concrete improvements over the inspected
sources. We do not claim to have established worldwide priority.

The factorial-free binomial CRT compiler and its two complementary
prime-power lemmas are inherited from the incoming factorial-free
submission~\cite{factorialfree}. Centered common remainders, the semantic
Pell anchor, signed factorization, and the integral balancing operations
are also retained and credited~\cite{cofinal,normalized}. The signed-digit
idea is adapted from the Shared Factorial Error submission~\cite{shared}.
The compact comparison and its native calibration come from the incoming
Lucas-correction and cofinal-witness submissions~\cite{lucas,cofinal}.

\section{The calibrated direct interface}
\label{sec:interface}

The direct route starts from the fixed numerical calculus $\calC$.
Its finitely described local rules have one conclusion and at most two
premises. The truth label is $0$ and the refutation label is $1$.
The retained local polynomial tests whether a row is a valid instance
and enforces two adjacent beta-code residues. A polynomially derived
local label sum $J$ bounds every label used by that row.

The logical endpoint is the composition
\[
\calC\text{ derives }1
\quad\longleftrightarrow\quad
\text{a finite ZF refutation exists}
\quad\longleftrightarrow\quad
\text{a finite ZFC refutation exists}.
\]
The named retained results are
\nolinkurl{KEndpoint.c222_iff_zf_refutation} and
\nolinkurl{ConstructibleChoice.zf_refutation_iff_zfc_refutation}.
Their pinned source files, syntax calibration, and ordinary bridge proofs
are included in the dependency tree~\cite{calibration}. Although the
same calculus is used elsewhere in proofs about machines, the present
direct polynomial compiles its numerical rules directly. It does not
encode a run of K215 and does not use Jones's universal expression.

Let $t$ denote the number of padded local rows. The actual CRT
representatives are $C$ and eight values
$A_a,A_b,A_c,A_d,A_S,A_{S'},A_u,A_v$. In the raw source several stored
coordinates are these values minus one. The literal auditors use those
shifts explicitly; semantic names in this report refer to actual values.
The derived label $J(\boldsymbol A)$ is evaluated from those
representatives. It is \emph{not} an extra independent CRT column.

For a bounded decoded local tuple $\boldsymbol w$, the retained
nonnegative polynomial $\LL$ has the first of the fixed majorants below;
the second is a separate bound on the globally encoded derived label:
\begin{equation}\label{eq:majorants}
0\le\LL(i,\boldsymbol w),J(\boldsymbol w)<(2+t+x+y+F)^E,
\qquad E=8{,}388{,}608,
\end{equation}
and $J(\boldsymbol A)\le H^\ell$ with $\ell=4{,}194{,}304$.
where $H=2+C+\sum A_j$ for nonnegative representatives. Its precise
arithmetic graph and the proofs of these generous bounds are retained.
Changing a range from $\le t$ to $<t$ preserves these upper bounds.

\subsection{Pell requests and signed bases}

For an integer center $A$, define Pell polynomials by
\[
X_0=1,\quad Y_0=0,\qquad
X_{k+1}=AX_k+(A^2-1)Y_k,\quad
Y_{k+1}=X_k+AY_k,
\]
with $X_{-k}=X_k$ and $Y_{-k}=-Y_k$. The retained square frame,
pooled index residues, and one semantic anchor force these exact values
at the requested signed indices. They give a frame $W\ge258$ that
strictly bounds all actual base, index, and output atoms and satisfies
$|A|>W^W$. The anchor uses the \emph{encoded} derived label
$J(\boldsymbol A)$; the small decoded local label is a different value.
No step here substitutes the latter for the former.

For a signed base $n$ the power congruence uses
\[
T_n=2An-n^2-1.
\]
The unit property of $n$ modulo $T_n$, the Pell recurrence, and the
absolute error bound imply the following inherited consequence:
\begin{equation}\label{eq:signed-power}
|n|\ge2\quad\Longrightarrow\quad
\text{a certified request }(n,k,m)\text{ forces }k\ge0, m=n^k.
\end{equation}
For negative $k$ the congruence would force $mn^{-k}=1$, impossible
for $|n|\ge2$. For nonnegative $k$ it forces $m=n^k$, since the
possible error is smaller in absolute value than $|T_n|$.
This statement allows negative bases. That feature is essential below.
Variable powers in the semantic discussion are conclusions of these
ordinary quadratic certificates, not operations in the final polynomial.

\section{Strict horizons remove a whole power certificate}
\label{sec:strict}

In the matched parent, four local schema parameters and the derived
label have bounds $a,b,c,d,J\le t$. Their binomial ranges use lower
index $t+1$ and denominator $q^{t+1}$. The index column already has
bound $C_{\rm local}<t$ and denominator $q^t$.

Choose instead
\begin{equation}\label{eq:strict}
t=16+h^2,\qquad a,b,c,d,J<t.
\end{equation}
Completeness is preserved: a finite proof, its finitely many row
parameters, and its derived labels can all be fixed \emph{before}
choosing a larger horizon. The fixed truth-padding row also has bounded
parameters and derived label. Taking a sufficiently large allowed
horizon dominates this finite set; it does not create a dependence on
the later CRT representatives or Pell witnesses.

The five range requests now use $\binom{A}{t}$, with the appropriate
actual upper representative $A$. They share the existing $q^t$
denominator and its normalization ratio. Hence we remove the request
for $q^{t+1}$, its one-step Pell addition graph, eight private source
coordinates, and the duplicate denominator-ratio equation. Some of those
coordinates would already disappear in monic elimination; the final
count is obtained from the emitted graph, not from adding raw removals.

The surviving 17 requests are exactly as follows. Identical indices
share their principal Pell pair.
\begin{center}\small
\begin{tabular}{lll}
\toprule
Base & Index or indices & Purpose\\
\midrule
$t$ &$t$ &$F=t^t$\\
$B$ &$x$ &$N=B^x$\\
$BH$ &$x$ &$q=(BH)^x$\\
$q+1$ &$F+t$ &exact binomial carrier\\
$q$ &$t,\ y+1,\ x$ &three denominator groups\\
$q+1$ &$A_a,A_b,A_c,A_d,A_S,A_{S'},A_u,A_v,C$
 &nine independent upper representatives\\
$q+1$ &$J(\boldsymbol A)$ &derived-label range and anchor\\
\bottomrule
\end{tabular}
\end{center}
There remain fourteen principal indexed Pell pairs and five base
moduli. The frame deletes the obsolete output and the obsolete $t+1$
index. Since the first base is now the existing index $t$, one duplicate
square atom is also removed. Every required atom remains bounded.
The new constructive frame is rebuilt; no equality with an old frame
witness is claimed.

\section{A self-power and a linear CRT radius}
\label{sec:growth}

The selected arithmetic choices are
\begin{equation}\label{eq:selected}
F=t^t,\qquad R_{\rm CRT}=F,\qquad
B=\binom{F+t}{t},\qquad y=F(2v+1)^2.
\end{equation}
Only the first equality needs a power certificate, at the already
present index $t$. The other displayed expressions describe the
mathematical meanings of literal quadratic registers and binomial digits.

\subsection{Label multiplicities at the self-power bound}

For each nonnegative label put $W_j=1+jF$. The commutative proof
certificate eventually gives
\begin{equation}\label{eq:balance}
\prod_{i<t}W_{v_i}=(1+F)\prod_{i<t}W_{p_i}W_{q_i}.
\end{equation}
Unused premise slots have label zero and contribute one.

\begin{lemma}[Self-power separation]\label{lem:separation}
Let $t\ge2$, $F=t^t$, and suppose all labels in
\eqref{eq:balance} are at most $t$. Then for every positive label $j$,
its number of conclusion occurrences minus its number of premise
occurrences is one for $j=1$ and zero otherwise.
\end{lemma}
\begin{proof}
Discard all factors of label zero. If the remaining factor counts
differ, their minimum $r$ is at most $t$, because the conclusion side
has at most $t$ factors. The larger product is at least
$(F+1)^{r+1}$ and the smaller at most $t^r(F+1)^r$. But
$F+1>t^t\ge t^r$, a contradiction. Thus both counts are the same
$r\le t$.

For a product $U(z)=\prod_{i=1}^r(1+a_i z)$ with
$1\le a_i\le t$, every coefficient $u_j$ is positive and
\[
1\le u_j\le\binom rj t^j\le r^{r-j}t^j\le t^r\le F.
\]
Therefore corresponding coefficients of the two product polynomials
differ by at most $F-1$. This difference bound remains strict even
when an individual coefficient equals $F$. If a nonzero difference
polynomial vanished at $F$, take its lowest nonzero coefficient,
divide the value by the corresponding power of $F$, and reduce
modulo $F$. This would force a nonzero integer of absolute value
less than $F$ to be divisible by $F$, a contradiction. Consequently
the products agree in $\Z[z]$. Their distinct factors $1+jz$ have
distinct roots $-1/j$, so unique factorization recovers exactly the
claimed multiplicities.
\end{proof}

The lemma is deliberately stated for the slightly larger range
$j\le t$; the new strict ranges are more than sufficient. It avoids
the incorrect claim that every individual coefficient is strictly
smaller than $t^t$.

\begin{lemma}[Single-conclusion flow]
If the factors in \eqref{eq:balance} come from valid single-conclusion
Horn rules and truth is derivable, then label $1$ is derivable.
\end{lemma}
\begin{proof}
Take the closure of truth under the finitely many listed rule
instances. If label $1$ is outside, any rule with conclusion outside
has at least one premise outside. Summing conclusions minus premises
over outside labels is at most zero. Lemma~\ref{lem:separation} makes
that sum one, a contradiction. This argument permits repeated
premises and cycles; a cycle cannot supply a positive net conclusion.
\end{proof}

\subsection{The radius need only be linear in the carrier}

\begin{lemma}\label{lem:lcm-bound}
For $t\ge1$, $\Lambda_t=\lcm(1,\ldots,t)$ satisfies
$\Lambda_t\le4^{t-1}$. For $t\ge16$,
\[
\Lambda_t^2\le16^{t-1}<16^t\le t^t=F.
\]
\end{lemma}
\begin{proof}
Let $k=\lceil n/2\rceil$. A prime power present in $\Lambda_n$
and missing from $\Lambda_k$ is greater than $k$ and at most $n$;
there is at most one additional power of that prime. Its contribution
to $\binom n{\lfloor n/2\rfloor}$ has valuation at least one,
by the floor-function formula. Thus
\[
\Lambda_n\mid\Lambda_k\binom n{\lfloor n/2\rfloor}.
\]
For $n\ge2$ the central binomial coefficient is at most $2^{n-1}$.
For odd $n$ this follows from its equal central neighbor; for even
$n=2m$ use $\binom{2m}m=2\binom{2m-1}{m-1}$ and the odd case.
Induction gives
$\Lambda_n\le4^{k-1}2^{n-1}\le4^{n-1}$, with initial value
$\Lambda_1=1$. Squaring and using $t\ge16$ proves the final chain.
\end{proof}

The matched parent used the radius $F^2$. Its generalized binomial
CRT theorem needs only $R_{\rm CRT}\ge\Lambda_t^2$. The lemma
therefore licenses $R_{\rm CRT}=F$, deleting a multiplication without
adding a certificate for a factorial or for an lcm.

\subsection{The two binomial prime-power facts}

For clarity we include the inherited facts that make the radius
argument sufficient. They have different roles in the two directions.

\begin{lemma}[Every prime has a compatible row]\label{lem:all-primes}
For $R,t\ge1$, put $B=\binom{R+t}{t}$. If $p^e\Vert B$, then
$p^e\mid R+i$ for some $1\le i\le t$.
\end{lemma}
\begin{proof}
Let $v=\max_i v_p(R+i)$. The factorial valuation for $B$ is a sum
of zero-one terms
\[
e=\sum_{a\ge1}\left(
\left\lfloor\frac{R+t}{p^a}\right\rfloor-
\left\lfloor\frac R{p^a}\right\rfloor-
\left\lfloor\frac t{p^a}\right\rfloor\right).
\]
For $a>v$ the interval contains no multiple of $p^a$. Such a power
also exceeds $t$, so that summand is zero. Therefore $e\le v$,
and a row attaining the maximum is compatible.
\end{proof}

\begin{lemma}[Every row has a large selector]\label{lem:selector}
If $R\ge\Lambda_t^2$, then for each $1\le i\le t$ there is
$p^e\Vert\binom{R+t}{t}$ with $p^e\mid R+i$ and $p^e>t$.
\end{lemma}
\begin{proof}
First observe the classical divisibility
\[
t\binom{t-1}{i-1}\mid\Lambda_t.
\]
Indeed its valuation at level $p^a$ is
$\lfloor t/p^a\rfloor-\lfloor(i-1)/p^a\rfloor-
\lfloor(t-i)/p^a\rfloor$, which is zero or one and vanishes for
$p^a>t$. This gives valuation at most $\lfloor\log_p t\rfloor$.

Since $R+i>\Lambda_t^2$, some prime has
$v=v_p(R+i)>2a$, where $a=v_p(\Lambda_t)$. Then $p^v>t$;
for $j\ne i$, $v_p(R+j)=v_p(j-i)$. Writing
$d=v_p(t\binom{t-1}{i-1})\le a$, the full valuation in $B$ is
$e=v-d$. Thus $a<e\le v$, so $p^e>t$ and $p^e\mid R+i$.
\end{proof}

\begin{lemma}[Binomial range]\label{lem:range}
If $A\in\N$, $k,e\ge1$, and $p^e\mid\binom Ak$, then
$A\bmod p^e<k$, with $\binom Ak=0$ when $A<k$.
\end{lemma}
\begin{proof}
The case $A<k$ is immediate. Otherwise, if no factor of
$A-k+1,\ldots,A$ is divisible by $p^e$, all valuation terms from
level $e$ onward are zero. The remaining zero-one terms total at
most $e-1$, a contradiction. Hence $A-r$ is divisible by $p^e$
for some $0\le r<k$, and the ordinary residue is less than $k$.
\end{proof}

For soundness, the constraints
\[
B\mid C+R+1,\qquad B^x\mid\binom Ct,
\qquad B^x\mid\binom{A_j}{b_j+1},\qquad
B^x\mid\LL(C,\boldsymbol A)
\]
select each row using Lemma~\ref{lem:selector}. Modulo its
$p^{xe}$, Lemma~\ref{lem:range} decodes bounded values. The raw
index equation and $p^e>t$ identify the selected row exactly.
If $0\le\LL<2^x\le p^{xe}$, its divisibility forces zero.

For completeness, use Lemma~\ref{lem:all-primes} at \emph{every}
prime of $B$. Prescribe a valid compatible row and its tuple modulo
$p^{2xe}$. CRT combines these conditions modulo $B^{2x}$. In a range
with lower index $k\le x$, its falling factorial has a factor divisible
by $p^{2xe}$. Division by $k!$ loses at most $v_p(k!)<x\le xe$,
so $p^{xe}$ still divides the binomial. This supplies all ranges and
the local congruence. A selector for each row alone would not prove
this converse; the all-primes lemma is necessary.

\section{An affine radix, with signs recovered in order}
\label{sec:radix}

Replace the parent's odd-square radix by
\begin{equation}\label{eq:radix}
q=4z+1,\qquad z\in\Z.
\end{equation}
This permits a linear elimination after all other transformations.
The residue class excludes $q=0,-1,-2$, so $|q+1|\ge2$ before
positivity is known. The weaker affine odd class $2z+1$ would admit
$q=-1$, whose numerator base is zero. The Shared Factorial Error
construction excludes unit radices using a separate factorial equation;
that equation is absent here. The modulus-four choice is what permits
the transfer of its signed-digit idea to this graph.

\subsection{Bootstrap before decoding the carrier}

The known positive base $t\ge16$ first gives $F=t^t>2^{32}$.
The odd-square beta multiplier gives $y\ge F>0$. All numerator
powers have base $q+1$ with absolute value at least two, so their
actual upper indices are nonnegative. In particular $C,A_j\ge0$
and $H\ge2$. The raw index equation is
\[
Bh_B=C+F+1.
\]
It excludes $B=0$, hence $|BH|\ge2$. The certified request
$q=(BH)^x$ now forces $x\ge0$, even if $B$ is provisionally negative.
The endpoints are
\begin{equation}\label{eq:endpoints}
x=F+1+(y+1)u_0,\qquad
x=1+[1+(t+1)y]u_t.
\end{equation}
The last equation excludes $x=0$. If $u_0\le-1$ the first gives
$x\le F-y\le0$, impossible. The last endpoint then has $u_t\ge1$.
Consequently
\begin{equation}\label{eq:amplifier}
x\ge F+1,\qquad x\ge2+(t+1)y\ge2+(t+1)F.
\end{equation}
With $Q=|q|$ we obtain
\[
Q=|BH|^x\ge2^x,\qquad |B|\le|BH|<Q/2.
\]
In particular $|q|\ge2$, so its denominator requests are genuine
signed powers. No claim $B>0$ or $q>0$ was used.

\subsection{Recover the exact binomial at either sign of the radix}

Put $n=F+t$ and $c_j=\binom nj$. For $0\le j\le t$,
\[
c_j\le(F+t)^t\le(F^2+F)^t\le2^{t(F-1)}.
\]
The last elementary inequality holds for $F\ge7$; it is much weaker
than the actual size of $F$. Together with \eqref{eq:amplifier} it gives
\[
\frac{\max_{j\le t}c_j}{Q-1}<\frac1{32},\qquad c_t<Q/32.
\]
The signed truncated polynomial
\[
r_* = \sum_{j=0}^{t-1}c_jq^j
\]
is congruent to $(q+1)^n$ modulo $q^t$ and satisfies
$|r_*|<Q^t/32$, regardless of the sign of $q$.

The common denominator is $M=q^x$. Every range or exact digit has a
denominator $D$ and integral ratio $v_D$ with $Dv_D=M$. A scaled row
has the form
\[
v_DU=MK+\rho.
\]
As $M\ne0$, also $v_D\ne0$; the row shows $v_D\mid\rho$.
Thus $r=\rho/v_D$ is an integer and $U=DK+r$. The common norm is
\begin{equation}\label{eq:norm}
M^2=1+4\sum_{j=1}^{11}\rho_j^2+\alpha^2+\beta^2+\gamma(\gamma+1).
\end{equation}
Every last term is nonnegative for integer witnesses. It follows that
$|r|<|D|/2$. In particular the carrier digit's remainder and $r_*$
are congruent modulo $q^t$ and differ in absolute value by less than
$(1/2+1/32)|q^t|$. They are equal.

Subtracting the low polynomial from $(q+1)^n$ and dividing by $q^t$
uses the literal carrier-digit quotient $K=B+q h_{\rm high}$:
\[
(q+1)^{F+t}=q^t(B+q h_{\rm high})+r.
\]
It therefore gives $B\equiv c_t\pmod q$. But
$|B-c_t|<Q/2+Q/32<Q$, so
\[
B=\binom{F+t}{t}>0,\qquad q=(BH)^x>0.
\]
The signs are consequences of the exact digit, rather than premises
of its decoding. The remaining arithmetic may now use $B\ge2$ and
$N=B^x>0$.

\subsection{The remaining digits, strict derived range, and proof flow}

For ordinary representatives $a\ge0$ and lower indices $j\le k\le x$,
$\binom aj\le\max(1,a)^k\le H^x$. For the derived label use
\eqref{eq:majorants} and $\ell t<x$ to obtain the same bound.
Since $q=B^xH^x$, all remaining genuine low remainders are less than
their denominator divided by $32$. The norm selects them uniquely.
There are now three denominator groups: $q^t$, $q^{y+1}$, and $q^x$.
The first serves the carrier, the index range, and all five strict ranges.

The derived range gives, at a selected prime power,
$J(\boldsymbol A)\equiv r\pmod{p^{xe}}$ with $0\le r<t$.
Column congruences give $J(\boldsymbol w)\equiv r$ modulo the same
power. The inherited local bound gives
$0\le J(\boldsymbol w)<2^x\le p^{xe}$, hence
$J(\boldsymbol w)=r<t$. This is the required justification for a
derived column; treating it as an independent representative would
not suffice.

The bounds $y\ge F$, \eqref{eq:amplifier}, and $t\ge16$ give
$2+t+x+y+F<2x$ and $x>2^{32}$. The elementary dyadic estimate
$2^{23}\log_2(2x)<x$ therefore yields
\[
0\le\LL<(2x)^{2^{23}}<2^x.
\]
Lemma~\ref{lem:lcm-bound} licenses radius $F$; the binomial CRT
argument recovers every bounded valid local row.

The algebraic beta choice permits $y=F$. In that case the first
endpoint has residue zero modulo $y+1$, so its initial accumulator
is zero. Each local update
\[
S_{i+1}W_{v_i}=S_iW_{p_i}W_{q_i}
\]
propagates zero because all weights are positive. The last endpoint
has residue one, a contradiction. Thus an actual zero has $y>F$,
initial accumulator $F+1$, and final accumulator one. Multiplying
the updates gives \eqref{eq:balance}; label separation and
single-conclusion flow give a finite $\calC$ refutation. The retained
logical calibration gives a finite ZFC refutation.

\section{Simultaneous constructive completeness}
\label{sec:complete}

Start with a finite calibrated $\calC$ refutation, expand it to a proof
tree, and fix its finite rule parameters and label data. The accumulator
and beta-quotient values will be chosen afterward. Choose a sufficiently large
\emph{even} $h$ so that $t=16+h^2$ strictly exceeds the tree size,
all schema parameters, all local derived labels, and those of the
fixed truth-padding row. Set $F=t^t$ and form the pending-multiset
accumulators using weights $1+jF$. Truth padding gives exactly $t$
rows. This chooses the horizon before all arithmetic witnesses.

The parity choice is essential for this self-power:
\[
v_2(F)=t\,v_2(t)\ge t>v_2(t!).
\]
Let $O$ be the odd part of $t!$ and choose an arbitrarily large odd
multiple $w$ of $O$. Set $y=Fw^2$, large enough to exceed $F$ and
all pending accumulators. Then $y$ is an even multiple of $t!$.
The beta moduli $1+(i+1)y$, $0\le i\le t$, are pairwise coprime:
a common prime for two of them would divide their index difference,
then $y$, then one of these numbers, an impossibility. They are odd.
CRT supplies a common beta code $x$; adding their odd product lets
us choose $x$ even and large. The endpoint quotients and all local
beta quotients meet their required ranges. The odd $w$ is exactly
$2v+1$ for an integer $v$.

Define $B=\binom{F+t}{t}$ and $N=B^x$. The binomial CRT construction
supplies the nine representative classes modulo $N^2$. For each
$1\le i\le t$, $v_2(F)>\lfloor\log_2t\rfloor$ gives
$v_2(F+i)=v_2(i)$, so $B$ is odd. Thus each representative class
has freely selectable parity. Choose the dominant schema
representative $a$ odd in $[2N^2,4N^2)$ and the other eight
representatives positive even and at most $2N^2$.

The retained shell-label construction has
$J(\boldsymbol A)\equiv1+a+b\pmod2$ and
$J(\boldsymbol A)\ge2a^2$. Hence the semantic anchor is even
and dominates every other index. Indeed
$F+t\le B<N^2$, $x,y+1,t<N^2$, while the anchor is at least
$8N^4$. The derived strict range is preserved because evaluating
its polynomial respects the CRT congruences modulo $p^{2xe}$;
division by $t!$ loses fewer than $x$ powers of $p$.

Now $H=2+C+\sum A_j$ is odd. With $x$ even,
$q=(BH)^x$ is an odd square, hence $q\equiv1\pmod4$ and
$z=(q-1)/4$ is an integer. These canonical witnesses satisfy the
new affine radix. The true normalized remainders obey
$0\le\rho_j<M/32$, and
\[
4\sum_{j=1}^{11}\rho_j^2<\frac{44}{1024}M^2<M^2.
\]
The nonnegative integral gap $g=M^2-1-4\sum\rho_j^2$ has the
form $\alpha^2+\beta^2+\gamma(\gamma+1)$. To see this, represent
$4g+1$ by three squares; it is $1$ modulo $4$, so exactly one
square is odd. Writing the two even roots as $2\alpha,2\beta$
and the odd one as $2\gamma+1$ gives the claimed identity.

The inherited constructive frame and semantic-anchor theorem supply
fresh simultaneous Pell witnesses for the seventeen genuine power
requests. Removing a duplicate frame atom does not obstruct that
construction. All other certified rows now have their intended
values. Thus every finite ZFC refutation gives a zero of the changed
arithmetic source, and the exact integral transformations give a zero
of each final quartic.

Factorials occur here only to prove existence of a suitable multiple
in the converse. No unknown in the new certificate is required to
equal $t!$. In particular the change from $[2(t+1)]^t$ to $t^t$
must use the even cofinal choice above; copying an odd-$h$ choice
from an intermediate $(t+1)^t$ variant would be wrong.

\section{From the source to the literal equations}
\label{sec:finite}

The final source retains only ordinary quadratic residuals. The
signed Horn factorization uses the exact identity
\[
\left(\prod L_i^2\right)(U^2+V^2)
=\left(\prod L_i\right)^2(U^2+V^2)
\]
and an analogous identity for the guarded local rule. The literal
factorization certificate is independently checked. Then 198 monic
pivots are eliminated. A pivot $x-f(\boldsymbol y)=0$ permits
integral substitution in every remaining row, and its inverse
restores $x=f(\boldsymbol y)$; each chosen substitution preserves
degree at most two. Unused coordinates may be dropped. Recorded
integral translations give the 239-variable, 176-row source.

Both endpoints apply six integral coordinate shears and five
preliminary row operations. Their coefficients and order are in
the certificate. Each shear is invertible over the integers, and
a row replacement $R_i\leftarrow R_i+kR_j$ with $i\ne j$ has
the inverse $R_i\leftarrow R_i-kR_j$. Three affine projections
give the principal; four give the dense companion. The final
balancing uses respectively 33 and 35 further row operations.
One parent operation that used the deleted successor-power row
is omitted explicitly. The mappings are transported by original
raw-row identities, not by guessed row offsets.

These operations preserve the existence of an integer zero of
the simultaneous residual system. They need not preserve the
numerical value of a sum of squares away from its zero set. Over
the integers, $\sum R_j^2=0$ is equivalent to every $R_j=0$.
The collected expansions have nonzero quartic terms and all
coordinates occur, so their degrees and variable counts are exact.

\subsection{Finite verification and its limits}

The release separates the following checks from the ordinary proofs:
\begin{enumerate}
\item Exact reconstruction of the parent-to-source edit, with a
complete census of changed and deleted rows.
\item An independent raw-interface audit of all seventeen requests,
fourteen principal norms, five base moduli, 34 deduplicated frame
atoms, eleven cap rows, six anchor-assembly rows, all ten range
identities, and the common normalized remainder budget.
\item Independent Horn factorization, all 198 monic projections,
translations, endpoint maps, every collected quartic coefficient,
and coordinate liveness.
\item Fresh byte-identical reconstruction of the source and both
direct endpoints, and of the compact circuits and fixed coefficients.
\item Exact free-polynomial cut identities for every compact gate,
complete leading homogeneous forms, and deliberate corruption
controls. The leading forms prove the exact degrees; random
evaluations would not be sufficient.
\end{enumerate}
Normal Python and optimized Python runs are both supplied; the
audits do not rely on removable \texttt{assert} statements.
Negative controls alter actual certificates or gates and are required
to be rejected. Checksums fix the artifacts, but are not evidence
for the unbounded arithmetic lemmas.

The source calibration and unchanged Pell lemmas are inherited
mathematical dependencies. Their pinned files, ordinary proof
sources, and original notices are retained. The archive is not a
standalone Lean or Rocq project closing every imported theorem,
and no fresh full repository proof-assistant build was performed.
The principal new formalization work still needed is the strict
horizon, self-power, signed-radix, and free-quotient proof layer.

\section{A free quotient in the compact Pell encoding}
\label{sec:free-quotient}

This section gives two additional equations in ten strictly positive
integer variables. Their exact arithmetic costs and total degrees are
\(83/312\) and \(89/278\). The second improves the earlier \(89/288\)
and \(92/278\) points in cost and degree. The first improves the earlier
\(84/328\) point in both metrics. Other tradeoffs remain relevant.
These are ordinary proofs composed with the retained native/refutation
calibration, not new proof-assistant certificates or global optimality
claims.

\subsection{The change of arithmetic architecture}

The incoming constructions used \(Y=qs\) and
\(c=kY+\eta\). The factor \(q\) then propagated through
\(U=ic^2\), making the strong Pell part expensive in degree.
We instead take
\[
 Y=s,\qquad \lambda=q,
\]
and put \(\lambda\) into the two Pell bases. This retains the
needed divisibility while lowering the degree of \(c\).
The complete denominator is recovered only after exponential projection.
The positive shift \(T+\alpha\) is restored: it supplies a direct
early proof of \(b>c\) when \(Y\) is no longer divisible by \(q\).

\subsection{Complete polynomial specification}

The independent variables are
\[
 J,F,\alpha,i,T,w,s,n,\eta,\zeta\in\mathbb Z_{>0}.
\]
The fixed tuple \((\beta,K_0,M_C,M_F)\) is the complete directly
indexed BFS943 radix67 tuple retained with this component. Its largest
port has 592,884 bits. Every byte and its finite reconstruction are
included; a port name is not an extra unknown. In particular
\(\beta\ge31\) is odd and
\(K_0>0\), \(0<M_C<\beta<M_F<2\beta\).
Set
\[
 h=\beta J,\quad q=h+1,\quad C=q-F-\alpha,\quad Y=s.
\]
Choose either the plain chart or the corrected chart:
\[
\begin{array}{c|ccc}
 &W_0&X&D\\ \hline
 \text{plain}&qw&q^2w&q^7\\
 \text{corrected}&w&qw&6q^7
\end{array}
\]
Both use the following polynomial abbreviations:
\begin{align*}
 k&=\eta+\zeta,&c&=kY+\eta,\\
 A&=qY(X+1)+1,&\Delta&=A^2-1,&H&=4A-5,\\
 v&=Xq^2Y^2,&E&=1+v(v+1)k^2,\\
 R&=[h(q-F)+\alpha](q^2-1)+(M_C+qM_F)J,\\
 B&=(K_0+W_0)C-F,&U&=ic^2,\\
 a&=\Delta U^2E,&F_0&=1+a,&G&=EF_0,\\
 F_s&=1+4aF_0,&Q&=\Delta(F_s-1),\\
 b&=c(T+\alpha)-R,&V&=F_sb-c,\\
 K&=X(3c-2X)+2,&L&=bV.
\end{align*}
In the plain chart let
\[
 E_c=k-R-1-q^4s;
\]
in the corrected chart let
\[
 E_c=6(k-R-1-q^4s)-v(R+1)(R-1)(R+3).
\]
Finally define
\begin{align*}
 M&=hDH,\\
 Z&=hD\{H(n-LB)+K\}-H(DB+hE_c),\\
 P&=Q(Z^2-M^2GV^2)+(M^2G-Z^2).
\end{align*}
The requested single equation is \(P=0\). Every displayed abbreviation
is a fixed polynomial. There are no variable exponents, divisions,
factorials, or extra predicates in the emitted circuits.

\paragraph{Theorem.}
For either chart, ZFC proves
\[
 \bigl(\nexists J,F,\alpha,i,T,w,s,n,\eta,\zeta>0:\ P=0\bigr)
 \quad\longleftrightarrow\quad\operatorname{Con}(\mathrm{ZFC}).
\]

\subsection{Ordered soundness argument}

For proof notation only, write
\[
 t=B/h,\quad W=1+hL,\quad \kappa=K/H,
 \quad \sigma=Y/q^3,\quad z=Z/M.
\]
For the plain chart put \(e=k-R-1\), \(\rho=e/q^7\).
For the corrected chart put
\[
 \bar e=k-R-1-\frac{v(R+1)(R-1)(R+3)}6,
 \qquad\rho=\bar e/q^7.
\]
No integrality of these quotients is initially assumed. Identically,
\[
 P=(Q-1)Z^2-M^2G(QV^2-1),\qquad
 z=n+\kappa+\sigma-Wt-\rho.                       \tag{F1}
\]

\paragraph{Root extraction and early positive transport.}
At a positive zero, the identities
\[
 F_s=(1+2a)^2,\qquad Q=4\Delta^2U^2G,
 \qquad\gcd(E,F_0)=1
\]
and prime valuations give \(E=\tau^2\), \(F_0=f_0^2\),
\(G=g^2\), where \(g=\tau f_0>0\). The retained rational
adjacent-Pell descent gives a positive integer \(y\) such that
\[
 QV^2-(Q-1)y^2=1,\qquad z=\pm gy.
\]
With \(f=1+2a\), \(I=2Ug\), one has
\[
 f^2-\Delta I^2=1,\quad c^2\mid I,\quad Q=(\Delta I)^2.
\]
The signed folding theorem gives \(c=\psi_A(p)\) for positive odd
\(p\). Here \(\psi_A(j)\) denotes the Pell ordinate. Since
\(c\ge3\), \(p\ge3\); therefore
\[
 c\ge4A^2-1>4q^4.
\]
The unconditional mask estimate is
\(R<q^4+(\alpha+2)q^2\). Thus
\[
 b-c>\alpha(c-q^2)-q^4-2q^2
       \ge c-q^4-3q^2>0.
\]
This argument includes \(Y=1\) and does not assume a sign of \(C\).
The auxiliary Pell index is odd; because
\(V>c(f^2-1)>4Q\), its first and third values are excluded. Hence
\[
 V>8Q^2,\qquad b>4g,\qquad gy<\frac{3}{8}bV<\frac L2.
                                                               \tag{F2}
\]

For the plain chart, \(\rho<b/q^7<L/2\). If \(B\le-1\),
then \(Wt<-L\), and (F1)--(F2) contradict \(n>0\).
For the corrected chart the same proof applies whenever
\((R+1)(R-1)(R+3)\ge0\). Its remaining cases are \(R\le-4\)
and \(R=0\). The exact identity
\[
 R=(hC+q\alpha)(q^2-1)+(M_C+qM_F)J
\]
shows that, for \(R=-r\le-4\),
\(C<0\), \(|C|>r/q^3\), and
\(Wt<-(1+w|C|)L\). The adverse correction is less than
\(ws^2r^3/(6q^4)\), whereas
\(L>(F_s-1)r^2\) and \(F_s-1>s^2\), so
\(w|C|L\) strictly exceeds it. For \(R=0\) the correction is
only \(ws^2/(2q^4)\) and the same margin suffices.
Therefore both charts force \(B\ge0\), hence \(C>0\).

The calibrated mask bound is now licensed:
\[
 3q<R<q^4-2.
\]
The normalized strong-rank and signed-index theorems, using
\(c^2\mid I\), \(c>2R\), \(c>2p\),
\(V\equiv-c\pmod f\) and \(V\equiv-R\pmod c\), give
\(p=R\). In particular \(R\) is odd.

\paragraph{Authenticate the first index.}
The first square gives \(k=2\psi_{1+2v}(n_0)\).
Comparison of the two Pell bases gives \(1\le n_0<R\).
Because \(H\equiv-1\pmod q\), clearing (F1) in the plain chart
gives \(q^4\mid e\). The first recurrence gives
\(k\equiv2n_0\pmod{4v}\), and here \(q^4\mid v\).
Consequently \(q^4\mid 2n_0-R-1\); its absolute value is less
than \(R<q^4\), so \(2n_0=R+1\).

The corrected chart has only \(q^3\mid v\). Let
\(r_0=(R+1)/2\). The exact Lucas expansion gives
\[
 2\psi_{1+2v}(m)=2m+8v\binom{m+1}{3}+32v^2Z_m,
 \qquad Z_m\in\mathbb Z_{\ge0}.
\]
Thus \(\bar e\) is integral. Clearing (F1) and cancelling the
literal factor 6 gives
\(q^7\mid\bar e-q^4s\), hence \(q^4\mid\bar e\).
For \(d=2n_0-R-1\) one first obtains \(q^3\mid d\), and then
\[
 6(\bar e-d)=vd\{(2n_0)^2+(2n_0)(R+1)+(R+1)^2-4\}
                    +192v^2Z_{n_0}.
\]
The right side is divisible by \(q^6\). For every prime
\(\ell\mid q\), its consequence
\(v_\ell(\bar e-d)\ge6v_\ell(q)-1\ge4v_\ell(q)\)
also covers \(\ell=2,3\). Hence \(q^4\mid d\), and the
same strict bound forces \(d=0\). Thus both charts give
\[
 n_0=(R+1)/2.                                             \tag{F3}
\]

\paragraph{Project before recovering the denominator.}
Reduction of the literal numerator modulo \(H\) gives
\(H\mid hK\) in the plain chart and \(H\mid6hK\) in the
corrected chart. Only the power of \(q\) was cancelled; no coprimality
of \(H\) and 3 is assumed.

Write \(R=2r+1\), \(r\ge48\). The retained ratio theorem is
uniform for every \(\lambda\ge1\), with
\(A=\lambda Y(X+1)+1\) and first base
\(1+2X\lambda^2Y^2\). Apply it at \(\lambda=q\) and use
(F3) and \(Y<c/k<Y+1\). It gives
\[
 Y>X^r/3,\qquad X^{r+1}<3Y(X+1)<H.
\]
The universal projection therefore yields
\(H\mid d_0h(X-2^R)(2^RX+1)\), with \(d_0=1\) or 6.
The entire product has absolute value below \(H\):
\[
 q^3h|(X-2^R)(2^RX+1)|
 <4\cdot4^rX^6+8\cdot16^rX^5<X^{r+1}.
\]
The last inequality holds for \(X\ge32,r\ge48\); after division
by \(X^{r+1}\), the corner value is
\(2^{-117}+2^{-25}<1\), and both terms decrease.
Consequently \(X=2^R\) and \(\kappa\) is integral.

Now \(q\mid X\) implies \(q=2^m\), \(m\ge5\), and
\(R>3q>6m\), so \(q^6\mid X\) and \(q^8\mid v\).
For the plain chart (F3) gives \(4v\mid e\); for the corrected
chart it gives \(\bar e=32v^2Z_{r_0}\). Thus \(\rho\) is a
nonnegative integer in both cases. Finally (F1) gives
\(Y/q^3-WB/h\in\mathbb Z\). Reduction modulo \(q^3\)
and \(h\) proves
\[
 q^3\mid Y,\qquad h\mid B,\qquad B/h\ge0.
\]
The retained parity, exact half-binomial rounding and fresh native-source
reconstruction now give a syntactic ZFC refutation.

\subsection{Constructive converse and exact degree}

A refutation supplies canonical outer data \(q^3\mid Y\),
\(X=2^R\), \(R\equiv3\pmod4\), \(R>3q\), and
\(B/h\ge0\). Set \(s=Y\), and take \(w=X/q^2\) or
\(w=X/q\) for the two charts. Use fresh bases at \(\lambda=q\),
and form \(c=\psi_A(R)\),
\(k=2\psi_{1+2v}((R+1)/2)\),
\(\eta=c-kY\), \(\zeta=k-\eta\).
The uniform ratio gives their strict positivity. A Pell matrix period
modulo \(c^2\tau\) supplies the seed \((f_0,I_0)\), and
\(i=I_0/(c^2\tau)>0\). Double the seed as above. Positive
auxiliary reconstruction gives \(T_*>\alpha\), so
\(T=T_*-\alpha>0\). Set
\[
 n=gy+(1+hL)t+\rho-\kappa-Y/q^3.
\]
This is positive because
\(g\ge f_0>(A-1)I_0>2Xc\), whereas
\(\kappa<Xc\) and \(Y/q^3\le Y<Xc\).
Its carrier equals \(gy\), proving \(P=0\).

For \(q_0=\beta J\), \(k_0=\eta+\zeta\), the complete leading
forms are
\begin{align*}
 [P_{\rm plain}]_{312}
   &=-2^{10}q_0^{140}w^{34}s^{78}i^{14}k_0^{44}(q_0-F)^2,\\
 [P_{\rm corrected}]_{278}
   &=-9\cdot2^{12}q_0^{106}w^{34}s^{78}i^{14}k_0^{44}(q_0-F)^2.
\end{align*}
Both have 135 terms. The highest part of \(b\) comes from \(-R\),
not from \(c(T+\alpha)\); retaining this change is essential.
The independent checker proves 28 exact cut identities per circuit,
covers every gate, verifies these entire leading forms, authenticates
all coefficient bytes, and rejects twenty deliberate actual corruptions.
Normal and optimized Python give identical mathematical receipts.
The complete unbounded proof and retained ordinary dependencies are
provided in the component's \texttt{PROOF.md} and source directory.


\section{A cofinal Pell experiment that did not improve the endpoint}
\label{sec:experiment}

We also tested replacing the exact growth power by a cofinal Pell
certificate. Put $d=6(t+1)$ and impose
\[
Y=d(3v+1),\qquad F=Y^2,\qquad
X^2-d(d+2)F=1.
\]
The nonzero ordinate of this Pell equation has an index at least $d$,
because $Y_n(d+1)\equiv n\pmod d$. It follows that $F$ exceeds
the factorial and label-growth budgets. A cofinal positive completion
can use $n=d10^t$, provided the chosen horizon has $3\nmid h$;
then $9\mid d$ and the Taylor congruence forces
$Y_n(d+1)/d\equiv1\pmod3$. Additional divisibility supplies a
compatible beta base. The full argument and independent recurrence
checks are in \nolinkurl{experiments/cofinal_pell/}.

This alternative is mathematically valid but was \emph{not} selected:
in the matched compilation trial before the strict-horizon change,
it cost two more variables than the shared exact-power route.
This comparison is not against the later 236-variable endpoint.
A first unqualified cofinal congruence claim also
failed at $t=16$; the preserved negative control detects that failure,
and the corrected proof restricts the cofinal horizon. Documenting
both facts avoids promoting a plausible compression idea before its
arithmetic and actual endpoint cost are checked.

\section{Use and incorporation}

For a directly inspectable small-coefficient consistency equation,
use \eqref{eq:principal}. For the smallest number of integer variables
among the new direct quartics, use the dense companion. For a compact
straight-line expression in ten positive variables, use the corrected
$89/278$ chart, retaining its coefficient files and its distinct
calibration. Replacing positive variables by $u+1$ changes the domain
to nonnegative integers; it does not make an unrestricted-integer
equation without additional work or a different construction.

From the archive root, the reproducibility entry points are:
\begin{lstlisting}
python3 -B verify.py
python3 -B -O verify.py
python3 -B verify.py --rebuild
cd report
pdflatex -interaction=nonstopmode -halt-on-error main.tex
pdflatex -interaction=nonstopmode -halt-on-error main.tex
\end{lstlisting}
The mathematical reconstruction uses only Python's standard library
and no network. LaTeX is needed only to rebuild the report. The
integration notes propose repository locations and explain the
retained source boundary; no repository write or publication is
part of this archive.

The proved progress is local and concrete: a full power certificate
can be removed by changing the witness horizon convention, and a
Pell scale can be moved to reduce the degree of a compact quotient
construction. Stronger global claims about the simplest consistency
equation would require a chosen complexity measure and further results.

\appendix
\section{Complete principal quadratic residuals}
\label{sec:residuals}

All variables $x_0,\ldots,x_{235}$ range over unrestricted integers.
The following 173 polynomials, together with \eqref{eq:principal},
specify the complete equation. No additional predicates or side
conditions are understood. Repeated subscripts mean ordinary
multiplication or a displayed square. The readable source is generated
from the literal JSON and checked against it by the release verifier.

% Generated from equation.json; do not edit independently.
\begingroup
\small
\setlength{\abovedisplayskip}{4pt}
\setlength{\belowdisplayskip}{4pt}
\[
\begin{aligned}
R_{0} &={}x_{0} - x_{0} x_{2} - x_{2} + x_{7}
\end{aligned}
\]
\[
\begin{aligned}
R_{1} &={}-x_{2} x_{5} - x_{3} + x_{6} x_{7}
\end{aligned}
\]
\[
\begin{aligned}
R_{2} &={}-1 + x_{1} - x_{6} x_{7}
\end{aligned}
\]
\[
\begin{aligned}
R_{3} &={}1 + x_{3} - x_{4} x_{18} + x_{16}
\end{aligned}
\]
\[
\begin{aligned}
R_{4} &={}-1 - x_{8}^{2} - 2 x_{8} x_{9} - x_{9}^{2} + x_{19}
\end{aligned}
\]
\[
\begin{aligned}
R_{5} &={}-10 x_{8} + x_{8}^{2} - 6 x_{8} x_{9} - 8 x_{8} x_{19} - 20 x_{9} - 11 x_{9}^{2} - 16 x_{9} x_{19} - 15 x_{19} - 4 x_{19}^{2} + x_{20}
\end{aligned}
\]
\[
\begin{aligned}
R_{6} &={}-2 + 2 x_{8} - x_{8}^{2} - 4 x_{8} x_{20} - 4 x_{20} - 4 x_{20}^{2} + x_{21}
\end{aligned}
\]
\[
\begin{aligned}
R_{7} &={}1 + 6 x_{8} - 36 x_{8}^{2} - 24 x_{8} x_{21} + 2 x_{21} - 4 x_{21}^{2} + x_{22}
\end{aligned}
\]
\[
\begin{aligned}
R_{8} &={}-2 - x_{9} - x_{9}^{2} - 2 x_{9} x_{10} - x_{10} - x_{10}^{2} + x_{23}
\end{aligned}
\]
\[
\begin{aligned}
R_{9} &={}-8 + 7 x_{8} - x_{8}^{2} + 4 x_{8} x_{10} - 4 x_{8} x_{23} - 4 x_{10} + 4 x_{23} + x_{24} - x_{46}
\end{aligned}
\]
\[
\begin{aligned}
R_{10} &={}-1 + x_{8} - x_{8}^{2} - 2 x_{8} x_{10} + x_{10} - x_{10}^{2} + x_{25}
\end{aligned}
\]
\[
\begin{aligned}
R_{11} &={}23 + 16 x_{8} - 8 x_{8} x_{9} - 8 x_{8} x_{10} - 16 x_{8} x_{19} - 8 x_{8} x_{25} + 19 x_{9} + 9 x_{9}^{2} - 36 x_{9} x_{10}
\\ &\phantom{={}}- 12 x_{9} x_{25} + 16 x_{10} + 8 x_{10}^{2} - 24 x_{10} x_{19} - 8 x_{19} x_{25} - x_{20} + x_{26} - x_{31}
\end{aligned}
\]
\[
\begin{aligned}
R_{12} &={}-4 x_{24}^{2} - 8 x_{24} x_{26} - 4 x_{26}^{2} + x_{27}
\end{aligned}
\]
\[
\begin{aligned}
R_{13} &={}-2 - 8 x_{8} + 8 x_{8} x_{23} + 8 x_{8} x_{25} - 8 x_{9} + 8 x_{9} x_{23} + 8 x_{9} x_{25} - 4 x_{19} + 8 x_{23} - 4 x_{23}^{2}
\\ &\phantom{={}}- 8 x_{23} x_{25} + 8 x_{25} - 4 x_{25}^{2} + x_{28}
\end{aligned}
\]
\[
\begin{aligned}
R_{14} &={}6 + 4 x_{8} + 32 x_{8} x_{9} + 16 x_{8} x_{19} + 16 x_{8} x_{25} + 8 x_{8} x_{28} + 24 x_{9} - 32 x_{9} x_{25} - 16 x_{9} x_{28}
\\ &\phantom{={}}+ 12 x_{19} - 16 x_{19} x_{25} - 8 x_{19} x_{28} - x_{20} + 8 x_{25} - 16 x_{25}^{2} - 16 x_{25} x_{28} + 4 x_{28} - 4 x_{28}^{2}
\\ &\phantom{={}}+ x_{29}
\end{aligned}
\]
\[
\begin{aligned}
R_{15} &={}22 - 5 x_{8} - x_{8}^{2} - 8 x_{8} x_{9} - 12 x_{8} x_{10} - 4 x_{8} x_{23} - 8 x_{9} - 16 x_{9} x_{10} - 16 x_{9} x_{23} - 20 x_{10}
\\ &\phantom{={}}- 20 x_{10}^{2} - 24 x_{10} x_{23} - 28 x_{23} - 4 x_{23}^{2} + x_{30}
\end{aligned}
\]
\[
\begin{aligned}
R_{16} &={}19 - 25 x_{8} - 7 x_{8}^{2} - 8 x_{8} x_{9} - 30 x_{8} x_{10} - 16 x_{8} x_{25} + 3 x_{9} + 5 x_{9}^{2} - 4 x_{9} x_{25} - 27 x_{10}
\\ &\phantom{={}}- 21 x_{10}^{2} - 24 x_{10} x_{25} - 6 x_{23} - 17 x_{25} - 4 x_{25}^{2} + x_{31}
\end{aligned}
\]
\[
\begin{aligned}
R_{17} &={}2 - 4 x_{30}^{2} - 8 x_{30} x_{31} - 4 x_{31}^{2} + x_{32}
\end{aligned}
\]
\[
\begin{aligned}
R_{18} &={}-19 x_{8} + 8 x_{8}^{2} + 8 x_{8} x_{19} + 8 x_{8} x_{23} - 46 x_{9} - 24 x_{9}^{2} + 24 x_{9} x_{19} + 24 x_{9} x_{23} - 2 x_{10}
\\ &\phantom{={}}+ 4 x_{19} - 4 x_{19}^{2} - 8 x_{19} x_{23} + 14 x_{23} - 4 x_{23}^{2} - x_{25} + x_{33}
\end{aligned}
\]
\[
\begin{aligned}
R_{19} &={}-3 - 6 x_{8} - 4 x_{10} - 2 x_{25} - 4 x_{33} - 4 x_{33}^{2} + x_{34}
\end{aligned}
\]
\[
\begin{aligned}
R_{20} &={}8 + 42 x_{8} + 7 x_{8}^{2} + 28 x_{8} x_{20} + 48 x_{9} - 32 x_{9}^{2} - 32 x_{9} x_{19} + 64 x_{9} x_{20} - 40 x_{19}^{2}
\\ &\phantom{={}}+ 32 x_{19} x_{20} + 14 x_{20} - 4 x_{20}^{2} + x_{35} + 2 x_{38}
\end{aligned}
\]
\[
\begin{aligned}
R_{21} &={}1 + 6 x_{8} - 36 x_{8}^{2} - 24 x_{8} x_{35} + 2 x_{35} - 4 x_{35}^{2} + x_{36}
\end{aligned}
\]
\[
\begin{aligned}
R_{22} &={}-x_{9}^{2} - 4 x_{9} x_{36} - 4 x_{36}^{2} + x_{37}
\end{aligned}
\]
\[
\begin{aligned}
R_{23} &={}-6 - 2 x_{8} - 4 x_{9} - 14 x_{19} - 16 x_{19}^{2} + x_{38}
\end{aligned}
\]
\[
\begin{aligned}
R_{24} &={}24 + 21 x_{9} - 21 x_{9}^{2} - 18 x_{9} x_{10} - 24 x_{9} x_{23} - 24 x_{9} x_{38} + 3 x_{10} - 9 x_{10}^{2} - 16 x_{10} x_{38}
\\ &\phantom{={}}+ 3 x_{23} - 12 x_{23}^{2} - 8 x_{23} x_{38} + 4 x_{38} - 4 x_{38}^{2} + x_{39} + 2 x_{46}
\end{aligned}
\]
\[
\begin{aligned}
R_{25} &={}16 - 10 x_{8} - 20 x_{8}^{2} - 16 x_{8} x_{10} - 24 x_{8} x_{25} - 24 x_{8} x_{39} - 12 x_{10} - 16 x_{10} x_{25} - 16 x_{10} x_{39}
\\ &\phantom{={}}- 14 x_{25} - 4 x_{25}^{2} - 8 x_{25} x_{39} + 4 x_{39} - 4 x_{39}^{2} + x_{40}
\end{aligned}
\]
\[
\begin{aligned}
R_{26} &={}-x_{10}^{2} - 2 x_{10} x_{11} - x_{11}^{2} + x_{41}
\end{aligned}
\]
\[
\begin{aligned}
R_{27} &={}2 + x_{8}^{2} - 4 x_{8} x_{41} + 2 x_{10} - 2 x_{10}^{2} - 8 x_{10} x_{41} - 2 x_{25} + 4 x_{41} - 4 x_{41}^{2} + x_{42}
\end{aligned}
\]
\[
\begin{aligned}
R_{28} &={}-31 - 16 x_{10} - 16 x_{10} x_{23} + 16 x_{10} x_{25} - 8 x_{23} + 12 x_{23}^{2} - 8 x_{23} x_{25} + 24 x_{25} - 4 x_{25}^{2} + x_{43}
\\ &\phantom{={}}- 4 x_{46}
\end{aligned}
\]
\[
\begin{aligned}
R_{29} &={}-4 x_{10} - 16 x_{25}^{2} - 32 x_{25} x_{41} - 16 x_{25} x_{42} - 16 x_{25} x_{43} - 4 x_{41} - 16 x_{41}^{2} - 16 x_{41} x_{42}
\\ &\phantom{={}}- 16 x_{41} x_{43} - 2 x_{42} - 4 x_{42}^{2} - 8 x_{42} x_{43} - 4 x_{43}^{2} + x_{44}
\end{aligned}
\]
\[
\begin{aligned}
R_{30} &={}-16 x_{8}^{2} + 8 x_{8} x_{10} + 32 x_{8} x_{25} + 16 x_{8} x_{28} + x_{10} - x_{10}^{2} - 8 x_{10} x_{25} - 4 x_{10} x_{28} - 16 x_{25}^{2}
\\ &\phantom{={}}- 16 x_{25} x_{28} - 4 x_{28}^{2} + x_{45}
\end{aligned}
\]
\[
\begin{aligned}
R_{31} &={}-1 - 8 x_{10} - 4 x_{10}^{2} + 8 x_{10} x_{23} + 8 x_{23} - 4 x_{23}^{2} + x_{46}
\end{aligned}
\]
\[
\begin{aligned}
R_{32} &={}-1 - 9 x_{10}^{2} + 24 x_{10} x_{23} + 12 x_{10} x_{46} - 16 x_{23}^{2} - 16 x_{23} x_{46} - 4 x_{46}^{2} + x_{47}
\end{aligned}
\]
\[
\begin{aligned}
R_{33} &={}1 + x_{9} - x_{9}^{2} - 4 x_{9} x_{10} - 4 x_{9} x_{47} - 4 x_{10}^{2} - 8 x_{10} x_{47} - 4 x_{47}^{2} + x_{48}
\end{aligned}
\]
\[
\begin{aligned}
R_{34} &={}-1 - 4 x_{48} - 4 x_{48}^{2} + x_{49}
\end{aligned}
\]
\[
\begin{aligned}
R_{35} &={}-1 - x_{10}^{2} - 8 x_{10} x_{48} - 4 x_{10} x_{49} - 16 x_{48}^{2} - 16 x_{48} x_{49} - 4 x_{49}^{2} + x_{50}
\end{aligned}
\]
\[
\begin{aligned}
R_{36} &={}2 + x_{9} - x_{9}^{2} - 4 x_{9} x_{10} - 4 x_{9} x_{50} - 4 x_{10}^{2} - 8 x_{10} x_{50} - 4 x_{50}^{2} + x_{51}
\end{aligned}
\]
\[
\begin{aligned}
R_{37} &={}-4 x_{45}^{2} - 8 x_{45} x_{51} - 4 x_{51}^{2} + x_{52}
\end{aligned}
\]
\[
\begin{aligned}
R_{38} &={}-9 + 7 x_{8} - x_{8}^{2} - 8 x_{8} x_{45} - 4 x_{8} x_{52} + 24 x_{45} - 16 x_{45}^{2} - 16 x_{45} x_{52} + 12 x_{52} - 4 x_{52}^{2}
\\ &\phantom{={}}+ x_{53}
\end{aligned}
\]
\[
\begin{aligned}
R_{39} &={}-1 - 4 x_{53} - 4 x_{53}^{2} + x_{54}
\end{aligned}
\]
\[
\begin{aligned}
R_{40} &={}-x_{10}^{2} - 8 x_{10} x_{53} - 4 x_{10} x_{54} - 16 x_{53}^{2} - 16 x_{53} x_{54} - 4 x_{54}^{2} + x_{55}
\end{aligned}
\]
\[
\begin{aligned}
R_{41} &={}-x_{8} - x_{8}^{2} - 4 x_{8} x_{10} - 4 x_{8} x_{55} - 4 x_{10} - 4 x_{10}^{2} - 8 x_{10} x_{55} - 4 x_{55} - 4 x_{55}^{2} + x_{56}
\end{aligned}
\]
\[
\begin{aligned}
R_{42} &={}-10 - 12 x_{56} - 4 x_{56}^{2} + x_{57}
\end{aligned}
\]
\[
\begin{aligned}
R_{43} &={}2 + 4 x_{10} + 16 x_{10} x_{56} + 8 x_{10} x_{57} - 4 x_{23} - 16 x_{23} x_{56} - 8 x_{23} x_{57} - x_{46} + 8 x_{56} - 16 x_{56}^{2}
\\ &\phantom{={}}- 16 x_{56} x_{57} + 4 x_{57} - 4 x_{57}^{2} + x_{58}
\end{aligned}
\]
\[
\begin{aligned}
R_{44} &={}x_{9} - x_{9}^{2} + 8 x_{9} x_{10} - 8 x_{9} x_{23} - 4 x_{9} x_{58} - 16 x_{10}^{2} + 32 x_{10} x_{23} + 16 x_{10} x_{58} - 16 x_{23}^{2}
\\ &\phantom{={}}- 16 x_{23} x_{58} - 4 x_{58}^{2} + x_{59}
\end{aligned}
\]
\[
\begin{aligned}
R_{45} &={}-10 - 12 x_{51} - 4 x_{51}^{2} + x_{60}
\end{aligned}
\]
\[
\begin{aligned}
R_{46} &={}1 - 16 x_{51}^{2} - 16 x_{51} x_{59} - 16 x_{51} x_{60} - 4 x_{59}^{2} - 8 x_{59} x_{60} - 4 x_{60}^{2} + x_{61}
\end{aligned}
\]
\[
\begin{aligned}
R_{47} &={}x_{10} - x_{10}^{2} - 8 x_{10} x_{59} - 4 x_{10} x_{61} - 16 x_{59}^{2} - 16 x_{59} x_{61} - 4 x_{61}^{2} + x_{62}
\end{aligned}
\]
\[
\begin{aligned}
R_{48} &={}-6 x_{8} - 3 x_{8}^{2} - 18 x_{8} x_{9} - 16 x_{9} - 27 x_{9}^{2} + 4 x_{9} x_{10} + 3 x_{10} + 3 x_{10}^{2} - 4 x_{10} x_{19} - 19 x_{19}
\\ &\phantom{={}}- 8 x_{19}^{2} + x_{20} + x_{63} - 4 x_{76}
\end{aligned}
\]
\[
\begin{aligned}
R_{49} &={}1 - 8 x_{8} - 4 x_{8} x_{10} + 8 x_{8} x_{63} - 16 x_{9} - 8 x_{9} x_{10} + 16 x_{9} x_{63} + x_{10} + x_{10}^{2} + 4 x_{10} x_{19}
\\ &\phantom{={}}+ 8 x_{19} - 8 x_{19} x_{63} + 3 x_{63} - 4 x_{63}^{2} + x_{64}
\end{aligned}
\]
\[
\begin{aligned}
R_{50} &={}-2 - 12 x_{8} - 16 x_{8} x_{10} + 16 x_{8} x_{62} + 16 x_{8} x_{64} - 24 x_{9} - 32 x_{9} x_{10} + 32 x_{9} x_{62} + 32 x_{9} x_{64}
\\ &\phantom{={}}+ 4 x_{10} + 4 x_{10}^{2} + 16 x_{10} x_{19} + 12 x_{19} - 16 x_{19} x_{62} - 16 x_{19} x_{64} - 4 x_{62}^{2} - 8 x_{62} x_{64} - 4 x_{63}
\\ &\phantom{={}}- 2 x_{64} - 4 x_{64}^{2} + x_{65}
\end{aligned}
\]
\[
\begin{aligned}
R_{51} &={}-2 - 4 x_{65} - 4 x_{65}^{2} + x_{66}
\end{aligned}
\]
\[
\begin{aligned}
R_{52} &={}-x_{3} x_{13} + x_{67}
\end{aligned}
\]
\[
\begin{aligned}
R_{53} &={}x_{12} - x_{13} - 2 x_{22} x_{67} - x_{68}
\end{aligned}
\]
\[
\begin{aligned}
R_{54} &={}x_{12} - x_{13} - 4 x_{24} x_{67} - 2 x_{27} x_{67} - x_{69}
\end{aligned}
\]
\[
\begin{aligned}
R_{55} &={}-2 x_{29} x_{67} + 2 x_{44} x_{67} - 2 x_{67} - x_{70} + x_{75}
\end{aligned}
\]
\[
\begin{aligned}
R_{56} &={}x_{12} - x_{13} - 4 x_{30} x_{67} - 2 x_{32} x_{67} - x_{71}
\end{aligned}
\]
\[
\begin{aligned}
R_{57} &={}x_{12} - x_{13} - 2 x_{34} x_{67} + 6 x_{67} - x_{72}
\end{aligned}
\]
\[
\begin{aligned}
R_{58} &={}x_{12} - x_{13} - 4 x_{36} x_{67} - 2 x_{37} x_{67} - x_{73}
\end{aligned}
\]
\[
\begin{aligned}
R_{59} &={}x_{12} - x_{13} - 2 x_{40} x_{67} - x_{74}
\end{aligned}
\]
\[
\begin{aligned}
R_{60} &={}x_{12} - x_{13} - 2 x_{44} x_{67} - 4 x_{67} - x_{75}
\end{aligned}
\]
\[
\begin{aligned}
R_{61} &={}-x_{8} x_{9} + x_{8} x_{10} + x_{9} + x_{9} x_{10} - x_{10} - x_{10}^{2} + x_{76}
\end{aligned}
\]
\[
\begin{aligned}
R_{62} &={}1 + 2 x_{11} + x_{11}^{2} - x_{76}^{2} + x_{77}
\end{aligned}
\]
\[
\begin{aligned}
R_{63} &={}-x_{77}^{2} + x_{78}
\end{aligned}
\]
\[
\begin{aligned}
R_{64} &={}-4 x_{65} x_{67} - 2 x_{66} x_{67} + 2 x_{67} + x_{79}
\end{aligned}
\]
\[
\begin{aligned}
R_{65} &={}-x_{12}^{2} + 2 x_{12} x_{13} + 2 x_{12} x_{79} - x_{13}^{2} - 2 x_{13} x_{79} - x_{79}^{2} + x_{80}
\end{aligned}
\]
\[
\begin{aligned}
R_{66} &={}-1 - x_{3} x_{9} + x_{81}
\end{aligned}
\]
\[
\begin{aligned}
R_{67} &={}-x_{13} x_{81} + x_{83} x_{84} + x_{85}
\end{aligned}
\]
\[
\begin{aligned}
R_{68} &={}-1 + x_{3} - x_{3} x_{8} + x_{82}
\end{aligned}
\]
\[
\begin{aligned}
R_{69} &={}-1 + 2 x_{3} + 2 x_{3} x_{8} + 4 x_{3} x_{9} - 2 x_{3} x_{19} + x_{83}
\end{aligned}
\]
\[
\begin{aligned}
R_{70} &={}-x_{12} x_{82} + x_{84}
\end{aligned}
\]
\[
\begin{aligned}
R_{71} &={}-6 x_{3} x_{13} + x_{12} x_{81} + 12 x_{13} - 8 x_{13} x_{81} - 4 x_{13} x_{82} - x_{13} x_{83} + x_{86}
\end{aligned}
\]
\[
\begin{aligned}
R_{72} &={}-x_{12} x_{68} + x_{13} x_{68} + x_{87}
\end{aligned}
\]
\[
\begin{aligned}
R_{73} &={}-x_{69} x_{87} + x_{88}
\end{aligned}
\]
\[
\begin{aligned}
R_{74} &={}-x_{70} x_{88} + x_{89}
\end{aligned}
\]
\[
\begin{aligned}
R_{75} &={}-x_{71} x_{89} + x_{90}
\end{aligned}
\]
\[
\begin{aligned}
R_{76} &={}-x_{72} x_{90} + x_{91}
\end{aligned}
\]
\[
\begin{aligned}
R_{77} &={}-x_{73} x_{91} + x_{92}
\end{aligned}
\]
\[
\begin{aligned}
R_{78} &={}-x_{74} x_{92} + x_{93}
\end{aligned}
\]
\[
\begin{aligned}
R_{79} &={}-x_{75} x_{93} + x_{94}
\end{aligned}
\]
\[
\begin{aligned}
R_{80} &={}-x_{85} x_{94} + x_{95}
\end{aligned}
\]
\[
\begin{aligned}
R_{81} &={}-x_{86} x_{95} + x_{96}
\end{aligned}
\]
\[
\begin{aligned}
R_{82} &={}-x_{96}^{2} + x_{97}
\end{aligned}
\]
\[
\begin{aligned}
R_{83} &={}-x_{2} - x_{2} x_{16} + x_{16} + x_{98}
\end{aligned}
\]
\[
\begin{aligned}
R_{84} &={}-x_{14} x_{98} + x_{99}
\end{aligned}
\]
\[
\begin{aligned}
R_{85} &={}-2 x_{1}^{2} + 2 x_{1} x_{12} + 2 x_{1} x_{13} + 2 x_{1} x_{99} + 2 x_{1} x_{100} - x_{12}^{2} - 2 x_{12} x_{99} - x_{13}^{2}
\\ &\phantom{={}}- 2 x_{13} x_{100} + x_{17} x_{102} - x_{78} x_{97} - x_{80} x_{97} - x_{99}^{2} - x_{100}^{2}
\end{aligned}
\]
\[
\begin{aligned}
R_{86} &={}-x_{2} x_{15} + x_{15} - x_{15} x_{98} + x_{100}
\end{aligned}
\]
\[
\begin{aligned}
R_{87} &={}-x_{8} - x_{9} - 4 x_{19} - 2 x_{22} - 4 x_{24} - 2 x_{27} - 2 x_{29} - 4 x_{30} - 2 x_{32} - 2 x_{34} - 4 x_{36} - 2 x_{37}
\\ &\phantom{={}}- 2 x_{40} - 2 x_{44} - 4 x_{65} - 2 x_{66} + x_{101}
\end{aligned}
\]
\[
\begin{aligned}
R_{88} &={}-x_{105} - 4 x_{105} x_{233} + x_{107}
\end{aligned}
\]
\[
\begin{aligned}
R_{89} &={}-x_{4} x_{128} + x_{103} x_{225} - x_{106} - x_{107} x_{128}
\end{aligned}
\]
\[
\begin{aligned}
R_{90} &={}x_{108} x_{225} - x_{109} x_{131} - x_{110}
\end{aligned}
\]
\[
\begin{aligned}
R_{91} &={}x_{111} x_{225} - x_{112} x_{131} - x_{113}
\end{aligned}
\]
\[
\begin{aligned}
R_{92} &={}x_{114} x_{225} - x_{115} x_{131} - x_{116}
\end{aligned}
\]
\[
\begin{aligned}
R_{93} &={}x_{117} x_{225} - x_{118} x_{131} - x_{119}
\end{aligned}
\]
\[
\begin{aligned}
R_{94} &={}x_{120} x_{226} - x_{122} x_{131} - x_{123}
\end{aligned}
\]
\[
\begin{aligned}
R_{95} &={}x_{124} x_{226} - x_{125} x_{131} - x_{126}
\end{aligned}
\]
\[
\begin{aligned}
R_{96} &={}-x_{17} x_{128} + x_{131}
\end{aligned}
\]
\[
\begin{aligned}
R_{97} &={}x_{127} - x_{129} x_{131} - x_{130}
\end{aligned}
\]
\[
\begin{aligned}
R_{98} &={}-x_{131} x_{133} + x_{132} - x_{134}
\end{aligned}
\]
\[
\begin{aligned}
R_{99} &={}-x_{131} x_{136} + x_{135} x_{225} - x_{137}
\end{aligned}
\]
\[
\begin{aligned}
R_{100} &={}-x_{131} x_{139} + x_{138} x_{225} - x_{140}
\end{aligned}
\]
\[
\begin{aligned}
R_{101} &={}x_{141} - 2 x_{141} x_{143} + x_{144}
\end{aligned}
\]
\[
\begin{aligned}
R_{102} &={}-x_{141} x_{144} - x_{144} + x_{145}
\end{aligned}
\]
\[
\begin{aligned}
R_{103} &={}2 x_{142} + x_{142}^{2} + x_{144}^{2} - x_{145}^{2}
\end{aligned}
\]
\[
\begin{aligned}
R_{104} &={}-x_{142} x_{147} - x_{147} + x_{148}
\end{aligned}
\]
\[
\begin{aligned}
R_{105} &={}2 x_{146} + x_{146}^{2} + x_{147}^{2} - x_{148}^{2}
\end{aligned}
\]
\[
\begin{aligned}
R_{106} &={}1 - 2 x_{0} + x_{0}^{2} - 2 x_{0} x_{142} + x_{149}
\end{aligned}
\]
\[
\begin{aligned}
R_{107} &={}-1 - x_{0} x_{147} + x_{3} + x_{142} x_{147} - x_{146} + x_{147} + x_{149} x_{150}
\end{aligned}
\]
\[
\begin{aligned}
R_{108} &={}1 - 2 x_{4} + x_{4}^{2} - 2 x_{4} x_{142} + x_{151}
\end{aligned}
\]
\[
\begin{aligned}
R_{109} &={}-x_{4} x_{194} + x_{17} + x_{151} x_{152} - x_{193} + x_{195}
\end{aligned}
\]
\[
\begin{aligned}
R_{110} &={}-1 + 2 x_{142} x_{209} - x_{153} + 2 x_{209} - x_{209}^{2}
\end{aligned}
\]
\[
\begin{aligned}
R_{111} &={}1 + x_{153} x_{154} - x_{193} - x_{194} x_{209} + x_{195} + 4 x_{233}
\end{aligned}
\]
\[
\begin{aligned}
R_{112} &={}-x_{142} x_{156} - x_{156} + x_{157}
\end{aligned}
\]
\[
\begin{aligned}
R_{113} &={}2 x_{155} + x_{155}^{2} + x_{156}^{2} - x_{157}^{2}
\end{aligned}
\]
\[
\begin{aligned}
R_{114} &={}1 - 4 x_{142} - 8 x_{142} x_{233} + x_{158} + 8 x_{233} + 16 x_{233}^{2}
\end{aligned}
\]
\[
\begin{aligned}
R_{115} &={}-1 + x_{103} - x_{155} - 2 x_{156} - 4 x_{156} x_{233} + x_{157} + x_{158} x_{159}
\end{aligned}
\]
\[
\begin{aligned}
R_{116} &={}-1 + x_{104} - 2 x_{142} x_{160} - x_{146} - x_{147} - 4 x_{147} x_{233} + x_{148} + x_{158} x_{160} + x_{160}
\\ &\phantom{={}}+ 8 x_{160} x_{233}
\end{aligned}
\]
\[
\begin{aligned}
R_{117} &={}1 + x_{142} - x_{142} x_{162} - x_{162} + x_{163}
\end{aligned}
\]
\[
\begin{aligned}
R_{118} &={}1 + 2 x_{161} + x_{161}^{2} - 2 x_{162} + x_{162}^{2} - x_{163}^{2}
\end{aligned}
\]
\[
\begin{aligned}
R_{119} &={}1 + x_{108} + x_{158} x_{164} - x_{161} - 2 x_{162} - 4 x_{162} x_{233} + x_{163} + 4 x_{233}
\end{aligned}
\]
\[
\begin{aligned}
R_{120} &={}-x_{142} x_{166} - x_{166} + x_{167}
\end{aligned}
\]
\[
\begin{aligned}
R_{121} &={}2 x_{165} + x_{165}^{2} + x_{166}^{2} - x_{167}^{2}
\end{aligned}
\]
\[
\begin{aligned}
R_{122} &={}-1 + x_{111} + x_{158} x_{168} - x_{165} - 2 x_{166} - 4 x_{166} x_{233} + x_{167}
\end{aligned}
\]
\[
\begin{aligned}
R_{123} &={}-x_{142} x_{170} - x_{170} + x_{171}
\end{aligned}
\]
\[
\begin{aligned}
R_{124} &={}2 x_{169} + x_{169}^{2} + x_{170}^{2} - x_{171}^{2}
\end{aligned}
\]
\[
\begin{aligned}
R_{125} &={}-1 + x_{114} + x_{158} x_{172} - x_{169} - 2 x_{170} - 4 x_{170} x_{233} + x_{171}
\end{aligned}
\]
\[
\begin{aligned}
R_{126} &={}-x_{142} x_{174} - x_{174} + x_{175}
\end{aligned}
\]
\[
\begin{aligned}
R_{127} &={}2 x_{173} + x_{173}^{2} + x_{174}^{2} - x_{175}^{2}
\end{aligned}
\]
\[
\begin{aligned}
R_{128} &={}-1 + x_{117} + x_{158} x_{176} - x_{173} - 2 x_{174} - 4 x_{174} x_{233} + x_{175}
\end{aligned}
\]
\[
\begin{aligned}
R_{129} &={}-x_{142} x_{178} - x_{178} + x_{179}
\end{aligned}
\]
\[
\begin{aligned}
R_{130} &={}2 x_{177} + x_{177}^{2} + x_{178}^{2} - x_{179}^{2}
\end{aligned}
\]
\[
\begin{aligned}
R_{131} &={}-1 + x_{120} + x_{158} x_{180} - x_{177} - 2 x_{178} - 4 x_{178} x_{233} + x_{179}
\end{aligned}
\]
\[
\begin{aligned}
R_{132} &={}-x_{142} x_{182} - x_{182} + x_{183}
\end{aligned}
\]
\[
\begin{aligned}
R_{133} &={}2 x_{181} + x_{181}^{2} + x_{182}^{2} - x_{183}^{2}
\end{aligned}
\]
\[
\begin{aligned}
R_{134} &={}-1 + x_{121} - 2 x_{142} x_{184} + x_{158} x_{184} - x_{181} - x_{182} - 4 x_{182} x_{233} + x_{183} + x_{184}
\\ &\phantom{={}}+ 8 x_{184} x_{233}
\end{aligned}
\]
\[
\begin{aligned}
R_{135} &={}-x_{142} x_{186} - x_{186} + x_{187}
\end{aligned}
\]
\[
\begin{aligned}
R_{136} &={}2 x_{185} + x_{185}^{2} + x_{186}^{2} - x_{187}^{2}
\end{aligned}
\]
\[
\begin{aligned}
R_{137} &={}-1 + x_{124} + x_{158} x_{188} - x_{185} - 2 x_{186} - 4 x_{186} x_{233} + x_{187}
\end{aligned}
\]
\[
\begin{aligned}
R_{138} &={}-x_{142} x_{190} - x_{190} + x_{191}
\end{aligned}
\]
\[
\begin{aligned}
R_{139} &={}2 x_{189} + x_{189}^{2} + x_{190}^{2} - x_{191}^{2}
\end{aligned}
\]
\[
\begin{aligned}
R_{140} &={}-1 + x_{127} + x_{158} x_{192} - x_{189} - 2 x_{190} - 4 x_{190} x_{233} + x_{191}
\end{aligned}
\]
\[
\begin{aligned}
R_{141} &={}1 - x_{142} x_{194} - x_{194} + x_{195}
\end{aligned}
\]
\[
\begin{aligned}
R_{142} &={}-1 + 2 x_{193} + x_{193}^{2} + x_{194}^{2} - 2 x_{195} - x_{195}^{2}
\end{aligned}
\]
\[
\begin{aligned}
R_{143} &={}x_{128} - 2 x_{142} x_{196} + x_{158} x_{196} - x_{193} - x_{194} - 4 x_{194} x_{233} + x_{195} + x_{196} + 8 x_{196} x_{233}
\end{aligned}
\]
\[
\begin{aligned}
R_{144} &={}-x_{142} x_{198} - x_{198} + x_{199}
\end{aligned}
\]
\[
\begin{aligned}
R_{145} &={}2 x_{197} + x_{197}^{2} + x_{198}^{2} - x_{199}^{2}
\end{aligned}
\]
\[
\begin{aligned}
R_{146} &={}-1 + x_{132} + x_{158} x_{200} - x_{197} - 2 x_{198} - 4 x_{198} x_{233} + x_{199}
\end{aligned}
\]
\[
\begin{aligned}
R_{147} &={}-x_{142} x_{202} - x_{202} + x_{203}
\end{aligned}
\]
\[
\begin{aligned}
R_{148} &={}2 x_{201} + x_{201}^{2} + x_{202}^{2} - x_{203}^{2}
\end{aligned}
\]
\[
\begin{aligned}
R_{149} &={}-1 + x_{135} + x_{158} x_{204} - x_{201} - 2 x_{202} - 4 x_{202} x_{233} + x_{203}
\end{aligned}
\]
\[
\begin{aligned}
R_{150} &={}-x_{142} x_{206} - x_{206} + x_{207}
\end{aligned}
\]
\[
\begin{aligned}
R_{151} &={}2 x_{205} + x_{205}^{2} + x_{206}^{2} - x_{207}^{2}
\end{aligned}
\]
\[
\begin{aligned}
R_{152} &={}-1 + x_{138} + x_{158} x_{208} - x_{205} - 2 x_{206} - 4 x_{206} x_{233} + x_{207}
\end{aligned}
\]
\[
\begin{aligned}
R_{153} &={}-7 - 2 x_{0}^{2} - 2 x_{0} x_{3} - x_{1}^{2} - x_{2}^{2} - 2 x_{3}^{2} - x_{4}^{2} + 2 x_{8} + 2 x_{8} x_{9} - x_{10}^{2} - x_{11}^{2}
\\ &\phantom{={}}- x_{12}^{2} - x_{13}^{2} - x_{14}^{2} - x_{15}^{2} - x_{16}^{2} - x_{17}^{2} - x_{19} - x_{101}^{2} - x_{103}^{2} - x_{104}^{2}
\\ &\phantom{={}}- x_{108}^{2} - x_{111}^{2} - x_{114}^{2} - x_{117}^{2} - x_{120}^{2} - x_{121}^{2} - x_{124}^{2} - x_{127}^{2} - x_{128}^{2}
\\ &\phantom{={}}- x_{132}^{2} - x_{135}^{2} - x_{138}^{2} + x_{141} - x_{209}^{2} - 24 x_{233} - 32 x_{233}^{2}
\end{aligned}
\]
\[
\begin{aligned}
R_{154} &={}-x_{4} - x_{4} x_{8} - x_{4} x_{9} - x_{4} x_{10} - x_{4} x_{11} - x_{4} x_{12} - x_{4} x_{13} - x_{4} x_{14} - x_{4} x_{15}
\\ &\phantom{={}}- x_{4} x_{16} + x_{209}
\end{aligned}
\]
\[
\begin{aligned}
R_{155} &={}-16 + x_{0} - x_{210}^{2}
\end{aligned}
\]
\[
\begin{aligned}
R_{156} &={}-x_{142} x_{212} - x_{212} + x_{221}
\end{aligned}
\]
\[
\begin{aligned}
R_{157} &={}-1 + x_{211}^{2} + x_{212}^{2} - x_{221}^{2}
\end{aligned}
\]
\[
\begin{aligned}
R_{158} &={}-x_{214} - x_{214} x_{215} + x_{222}
\end{aligned}
\]
\[
\begin{aligned}
R_{159} &={}-1 + x_{213}^{2} + x_{214}^{2} - x_{222}^{2}
\end{aligned}
\]
\[
\begin{aligned}
R_{160} &={}-x_{206}^{2} + x_{223}
\end{aligned}
\]
\[
\begin{aligned}
R_{161} &={}x_{212} - 2 x_{217} x_{223} - x_{223}
\end{aligned}
\]
\[
\begin{aligned}
R_{162} &={}-4 x_{212} - 8 x_{212} x_{216} + x_{215}
\end{aligned}
\]
\[
\begin{aligned}
R_{163} &={}-x_{142} - x_{211} x_{218} + x_{215}
\end{aligned}
\]
\[
\begin{aligned}
R_{164} &={}-1 - x_{205} - x_{211} x_{219} + x_{213}
\end{aligned}
\]
\[
\begin{aligned}
R_{165} &={}-x_{101} - 4 x_{206} x_{220} + x_{214}
\end{aligned}
\]
\[
\begin{aligned}
R_{166} &={}-2 + 2 x_{0}^{2} + 2 x_{0} x_{3} - 2 x_{0} x_{147} - 2 x_{0} x_{156} + x_{1}^{2} - 2 x_{1} x_{194} + x_{2}^{2} - 2 x_{2} x_{182} + x_{3}^{2}
\\ &\phantom{={}}- 2 x_{3} x_{156} + x_{8}^{2} - 2 x_{8} x_{162} + x_{9}^{2} - 2 x_{9} x_{166} + x_{10}^{2} - 2 x_{10} x_{170} + x_{11}^{2} - 2 x_{11} x_{174}
\\ &\phantom{={}}+ x_{12}^{2} - 2 x_{12} x_{178} + x_{13}^{2} - 2 x_{13} x_{186} + x_{14}^{2} - 2 x_{14} x_{190} + x_{15}^{2} - 2 x_{15} x_{198} + x_{16}^{2}
\\ &\phantom{={}}- 2 x_{16} x_{202} - 2 x_{101} x_{206} - x_{142} x_{224} + x_{146}^{2} + x_{155}^{2} + x_{161}^{2} + 2 x_{162} + x_{165}^{2}
\\ &\phantom{={}}+ x_{169}^{2} + x_{173}^{2} + x_{177}^{2} + x_{181}^{2} + x_{185}^{2} + x_{189}^{2} + x_{193}^{2} + x_{197}^{2} + x_{201}^{2}
\\ &\phantom{={}}+ x_{205}^{2} + 2 x_{206}^{2} - x_{227}^{2} - x_{228}^{2} - x_{229} - x_{229}^{2}
\end{aligned}
\]
\[
\begin{aligned}
R_{167} &={}-x_{104} x_{225} + x_{128}
\end{aligned}
\]
\[
\begin{aligned}
R_{168} &={}-x_{121} x_{226} + x_{128}
\end{aligned}
\]
\[
\begin{aligned}
R_{169} &={}-1 + x_{2} - x_{3} x_{234}
\end{aligned}
\]
\[
\begin{aligned}
R_{170} &={}-2 - x_{101}^{2} - x_{147}^{2} - x_{156}^{2} + 2 x_{162} - x_{162}^{2} - x_{166}^{2} - x_{170}^{2} - x_{174}^{2} - x_{178}^{2}
\\ &\phantom{={}}- x_{182}^{2} - x_{186}^{2} - x_{190}^{2} - x_{194}^{2} - x_{198}^{2} - x_{202}^{2} + x_{206}^{2} - x_{227}^{2} - x_{228}^{2}
\\ &\phantom{={}}- x_{229} - x_{229}^{2}
\end{aligned}
\]
\[
\begin{aligned}
R_{171} &={}-1 - 4 x_{106}^{2} - 4 x_{110}^{2} - 4 x_{113}^{2} - 4 x_{116}^{2} - 4 x_{119}^{2} - 4 x_{123}^{2} - 4 x_{126}^{2} + x_{128}^{2} - 4 x_{130}^{2}
\\ &\phantom{={}}- 4 x_{134}^{2} - 4 x_{137}^{2} - 4 x_{140}^{2} - x_{230}^{2} - x_{231}^{2} - x_{232} - x_{232}^{2}
\end{aligned}
\]
\[
\begin{aligned}
R_{172} &={}-1 + x_{234} - 4 x_{235} - 4 x_{235}^{2}
\end{aligned}
\]
\endgroup


\section{Artifact identities and comparison scope}

\begin{center}\small
\begin{tabular}{ll}
\toprule
Object & Relative path\\
\midrule
Principal residuals &\nolinkurl{direct/principal236/equation.json}\\
Principal full quartic &\nolinkurl{direct/principal236/expanded_polynomial.json}\\
Dense residuals &\nolinkurl{direct/dense235/equation.json}\\
Dense full quartic &\nolinkurl{direct/dense235/expanded_polynomial.json}\\
Strict arithmetic source &\nolinkurl{direct/source/raw_source.json}\\
Plain compact circuit &\nolinkurl{compact/equation_83_degree312.json}\\
Corrected compact circuit &\nolinkurl{compact/equation_89_degree278.json}\\
Exact fixed compact tuple &\nolinkurl{compact/upstream_lucas/parent/coefficients/}\\
Proof dependencies &\nolinkurl{sources/factorial_free/} and \nolinkurl{compact/}\\
Independent reviews and receipts &\nolinkurl{review/}\\
\bottomrule
\end{tabular}
\end{center}

The principal residual JSON has SHA-256
\begin{center}\small\ttfamily
4e619e431c9beb13967aff3e7e7db35c56951fb179c353c43ee1fe7601dc3cbe.
\end{center}
The dense residual JSON has SHA-256
\begin{center}\small\ttfamily
9988b2f76ef9f2adb4d7f6f9867134b1dc695905c0fc1aa189da753a745949d4.
\end{center}
The root source ledger records inspected ZIP identities and the repository
snapshot; \nolinkurl{MANIFEST.sha256} fixes the delivered files. The
comparison covers the linked repository directories and their incoming
submissions at that snapshot. It does not establish a universal or
historical optimum.

\begin{thebibliography}{99}
\raggedright
\bibitem{repo} Vladimir Reshetnikov and contributors,
\emph{ProveIt}, repository snapshot
\texttt{7a8da414376a80b90ecf799f37ceb26fb9ac674e},
\url{https://github.com/VladimirReshetnikov/ProveIt/tree/7a8da414376a80b90ecf799f37ceb26fb9ac674e}.
The inspected directories are \nolinkurl{Computability/HilbertTenthProblem},
\nolinkurl{Computability/DiophantineConsistency/con-zfc-diophantine-forms},
and \nolinkurl{docs/incoming}.

\bibitem{factorialfree} Incoming ProveIt submission,
\emph{Removing Factorial Authenticity: Binomial Prime-Power Certificates
for Diophantine Consistency}, 25 October 2026,
\nolinkurl{ConZFC_Factorial_Free_Binomial_CRT_2026-10-25.zip}.
Complete supplied source retained under \nolinkurl{sources/factorial_free/}.

\bibitem{cofinal} Incoming ProveIt submission,
\emph{Cofinal Witnesses and Centered Carriers}, 25 October 2026,
\nolinkurl{ConZFC_Cofinal_Witnesses_Centered_Carriers_2026-10-25.zip}.
Relevant compact and direct dependencies are retained in their source chains.

\bibitem{normalized} Incoming ProveIt submission,
\emph{Normalized Residues, Round 19}, 25 October 2026,
\nolinkurl{ConZFC_Normalized_Residues_Round19_2026-10-25.zip}.
Retained within the direct dependency tree.

\bibitem{shared} Incoming ProveIt submission,
\emph{Shared Factorial Error}, 25 October 2026,
\nolinkurl{ConZFC_Shared_Factorial_Error_2026-10-25.zip}.
Source of the signed-binomial-digit comparison used here.

\bibitem{lucas} Incoming ProveIt submission,
\emph{Lucas Corrections and Witness Elimination}, 25 October 2026,
\nolinkurl{ConZFC_Lucas_Corrections_and_Witness_Elimination_2026-10-25.zip}.
Retained under \nolinkurl{compact/upstream_lucas/}.

\bibitem{calibration} ProveIt logical calibration, pinned source commit
\texttt{08200eb10fde2db408c69dffd840d6920a569b1f}.
See \nolinkurl{KEndpoint.lean}, \nolinkurl{ConstructibleChoice.lean},
and \nolinkurl{CALIBRATION_PINS.json} in
\nolinkurl{sources/factorial_free/direct/parent/upstream/round17/architecture/proofs/inherited/calibration/}.
The pins identify theorem statements and source provenance; the present
release does not claim a new full build of their dependency graph.
\end{thebibliography}
\end{document}
